 % ! TeX root = ../mat351.tex
\section{Day 43: More Applications of the Fourier Transform (Apr.\ 1, 2026)}
What does \( e^{- i t \Delta} \) mean?

The answer is 
\[ e^{- i t \Delta} \colon f \to u = e^{-i t \Delta} f \]
such that
\[ \left\{\begin{NiceMatrix}\partial_{t} u = - i \Delta u \\ u(0, x) = f & x \in \mathbb{R}^{d}\end{NiceMatrix}\right.  \]
set \( u = U(t) f \). Then we want that
\[ \partial_{t} u = - i \Delta u, \quad \left. e^{-i t \Delta} f \right|_{t = 0} = f \]
So one answer is through the Fourier transform.
\[ \left( e^{-i t \Delta} f \right)(t, x) = \hat {\mathcal{F}}^{-1}_{\xi \to x} \left( e^{i t 4 \pi^{2} | \xi|^{2}} \hat f( \xi) \right), \]
where we map the \( \xi \) variable to \( x \) (which is what this somewhat enigmatic subscript means). From the definition we can check that
\begin{align*}
	\partial_{t} \hat {\mathcal{F}}^{-1} \left( e^{i t 4 \pi^{2} |\xi|^{2}} \hat f(\xi) \right) &= \hat {\mathcal{F}}^{-1} \left( \partial_{t}e^{i t 4 \pi^{2} |\xi|^{2}} \hat f(\xi) \right) \\
																																															&= \hat {\mathcal{F}}^{-1}\left( e^{i t 4 \pi^{2} |\xi|^{2}} i 2 \pi \xi \cdot 2 \pi \xi \hat f(\xi) \right) \\
																																															&= i \hat {\mathcal{F}}^{-1} \left( 2 \pi \xi \cdot 2 \pi \xi e^{i t 4 \pi^{2} |\xi|^{2}}  \hat f(\xi)\right) \\
	&= -i \Delta \hat {\mathcal{F}}^{-1} \left( e^{i t 4 \pi^{2} |\xi|^{2}}\hat f(\xi) \right)
\end{align*}
where this last line is because \( \hat {\mathcal{F}}(\nabla g) = 2 \pi i \xi \hat g(\xi) \) and likewise \( \hat {\mathcal{F}}(\nabla \cdot G) = 2 \pi i \xi \hat G(\xi) \).

For a general function \( a \), we can define
\( a(\Delta) \colon f \to a(\Delta) f \)
by doing
\[ a(\Delta) f = \hat {\mathcal{F}}\left( a(2 \pi i \xi) \hat f(\xi) \right) \]

Last class we proved Heisenberg's uncertainty principle. Namely, if \( f \colon \mathbb{R} \to \mathbb{C} \) has \( \| f \|_{L^{2}} = 1 \), and
\[ \bar x = \int_{\mathbb{R}} x^{2} |f|^{2} \, dx \quad \bar k = \int_{\mathbb{R}} \xi^{2} |\hat f|^{2} d \xi, \]
then
\[ \hat x \cdot \hat k \ge \frac{1}{16 \pi^{2}}. \]
We also get the same with
\[ \bar x = \int_{} (x - x_{0})^{2} |f|^{2} \, dx \quad \bar k = \int (\xi - \xi_{0})^{2} |\hat f|^{2} \, d \xi. \]
How? Well apply the original to \( f(x + x_{0}) \) and we can shift \( \bar x \). This won't change the value of \( \bar k \). For shifting \( \bar k \), we should modulate \( f \). Namely, we should take \( e^{- 2 \pi i \xi_{0} x} f(x + x_{0}) \), which should shift \( \bar k \). Fabio says to check this yourself.

But now one can ask when is \( \bar x \bar k = \frac{1}{16 \pi^{2}} \)? Well we proved the inequality originally via Cauchy-Schwarz, so
\[ \left| \int x f\,  f' \right| = \| x f\|_{L^{2}} \| f' \|_{L^{2}} \]
implies that for some \( c \)
\[ x f(x) = c f'(x), \]
and therefore we have that for some \( A, B \)
\[ f(x) = Ae^{B\frac{x^{2}}{2}}. \]
But which choices of \( A, B \) yield \( \| f \|_{L^{2}} = 1 \)? We want \( \operatorname{Re} B < 0 \) so that it decays, therefore by using our knowledge of the Gaussian integral, that
\[ \int_{\mathbb{R}} e^{- \pi x^{2}} \, dx = 1,  \]
we arrive at
\[ \int_{\mathbb{R}} e^{- B x^{2}} \, dx = \frac{\sqrt{\pi}}{\sqrt{2 \operatorname{Re} B}}  \]
which then yields
\[ |A|^{2} = \frac{\sqrt{2 \operatorname{Re}(B)}}{\sqrt{\pi}} \]
to get that
\[ 1 = \int_{} |A|^{2} e^{- 2 \operatorname{Re}(B) x^{2}}\, dx.  \]
Fabio says to check this as well because he's not sure of it.

We are going to solve
\[ \left\{\begin{NiceMatrix}- \Delta u = 0 \\ u(x, 0) = g(x) \\ u(x, b) = h(x)\end{NiceMatrix}\right.  \]
on the domain of
\[ 0 < z < b \quad x \in \mathbb{R}^{d} \]
and rediscover D'Alembert's formula for
\[ \left\{\begin{NiceMatrix}u_{t t} - c^{2} u_{x x} = 0 \\ u(0, x) = g(x) \\ u_{t}(0, x) = h(x).\end{NiceMatrix}\right.  \]
For this second question, take the Fourier transform in \( x \) of
\[ \left\{\begin{NiceMatrix}\hat u_{t t} + (2 \pi c \xi)^{2} \hat u(\xi) = 0 \\ \hat u(0, \xi) = \hat g(\xi) \\ \hat u_{t}(0, \xi) = \hat h(\xi)\end{NiceMatrix}\right.  \]
This PDE has turned into a second order ODE\@. Back when we did this a while ago we spoke about the ``operator method'' and said ``treat the Laplacian as a matrix''. This is just that.

To solve our ODE, the general solution is
\[ \hat u(t, \xi) = A(\xi) \cos(2 \pi c \xi t) + B(\xi) \sin(2 \pi c \xi t). \]
Imposing our initial conditions gives us that
\[ \hat u(0, \xi) = A(\xi) \implies A = \hat g. \]
Likewise we also have that
\[ \partial_{t} \hat u(0, \xi) = B(\xi) 2 \pi c \xi = \hat h(\xi) \implies B(\xi) = \frac{\hat h(\xi)}{2 \pi c \xi}. \]
This implies that
\[ \hat u(t, \xi) = \cos(2 \pi c \xi t) \hat g(\xi) + \frac{\sin(2 \pi c \xi t)}{2 \pi c \xi}\hat h(\xi). \]
This is what we saw before. Rewriting this gives us
\[ u(t, x) = \cos\left(c \frac{\partial_{x}}{i} t\right)g + \frac{\sin \left( c \frac{\partial x}{i} t \right)}{-i \partial_x} h. \]
This is by our earlier proposition on \( a(\nabla) \) and such. This rewriting is just a curiosity though, it's not necessary for the solution. We'll now show it's equivalent to D'Alembert. How would we show this? Well first we need to show
\begin{align*}
	\hat {\mathcal{F}}^{-1} \left( \cos(2 \pi c \xi t) \hat g(\xi) \right) &= \frac{1}{2} \left( g(x + ct) + g(x - ct) \right) \\
	\cos(2 \pi c \xi t) \hat g(\xi) &= \frac{1}{2} \hat {\mathcal{F}} \left( g(x + ct) + g(x - ct) \right) \\ 
																	&= \frac{1}{2} \hat g(\xi) e^{2 \pi i c t \xi} + \frac{1}{2} \hat g (\xi) e^{- 2 \pi i c t \xi} \\
																	&= \cos(2 \pi c t \xi) \hat g(\xi)
\end{align*}
the second part for \( \sin \) you really should do, but he's not going to himself. Also we should do the first problem ourselves if we have time.

How do we define the Fourier Transform of a not nice function? If it's not \( L^{2} \) or not Schwartz. If \( \psi \) is in the dual of the Schwartz functions---a bounded linear function on \( S \)---then we say that
\[ (\hat \psi, \phi) = (\psi, \hat \phi), \quad \forall \phi \in S. \]

Thanks Fabio for the classes.
